\documentclass[a4paper, openany, 12pt]{article}

%% Стандарт библиографии — GOST 7.1 (gost71u.bst).
\bibliographystyle{gost71u}

%% подключаем преамбулу: в ней содержится подключение всех необходимых пакетов
\input{include/preambule.tex}

\begin{document}
    %% титульник
    \input{include/title-page.tex}
    \newpage

    \setcounter{page}{2}
    \fancyfoot[c]{\thepage}
    %% надпись над верхним колонтитулом
    \fancyhead[L]{\input{include/work-title.tex}}
    \fancyhead[R]{}

    %% аннотация
    \input{parts/Annotation.tex}
    \newpage

    %% содержание (генерируется автоматически по заголовкам)
    \tableofcontents{}
    \newpage

    %% обозначения и сокращения
    \input{parts/Notation.tex}
    \newpage

    %% введение
    \input{parts/Chapter0.tex}
    \newpage

    %% Глава 1. Физическая постановка задачи мультипольного анализа
    \input{parts/Chapter1.tex}
    \newpage

    %% Глава 2. Генерация обучающей выборки. Признаковое представление и предобработка
    \input{parts/Chapter2.tex}
    \newpage

    %% Глава 3. Архитектуры моделей
    \input{parts/Chapter3.tex}
    \newpage

    %% Глава 4. Функции потерь и обучение
    \input{parts/Chapter4.tex}
    \newpage

    %% Глава 5. Метрики качества и наборы данных оценки
    \input{parts/Chapter5.tex}
    \newpage

    %% Глава 6. Эксперименты и результаты
    \input{parts/Chapter6.tex}
    \newpage

    %% Глава 7. Программный фреймворк mpinv
    \input{parts/Chapter7.tex}
    \newpage

    %% Заключение
    \input{parts/Conclusion.tex}
    \newpage

    %% Список использованных источников — BibTeX по ГОСТ 7.1
    \renewcommand{\refname}{Список использованных источников}
    \addcontentsline{toc}{section}{Список использованных источников}
    \bibliography{references}

    %% В зависимости от надобности подключаем раздел "Приложение"
    % \newpage
    % \input{parts/Appendix.tex}
\end{document}
